\documentclass[11pt]{amsart}
\usepackage{amsmath,amssymb}
\setlength{\emergencystretch}{3em}
\tolerance=600
\usepackage{tikz-cd}
\usepackage{tikz}
\usepackage[hidelinks]{hyperref}

\newtheorem{thm}{Theorem}[section]
\newtheorem{lem}[thm]{Lemma}
\newtheorem{cor}[thm]{Corollary}
\newtheorem{prop}[thm]{Proposition}
\theoremstyle{definition}
\newtheorem{defi}[thm]{Definition}
\newtheorem{rem}[thm]{Remark}
\newtheorem{ex}[thm]{Example}
\newtheorem{que}[thm]{Question}

\newcommand{\Set}{\mathbf{Set}}
\newcommand{\Fin}{\mathbf{FinSet}}
\newcommand{\Sub}{\mathrm{Sub}}
\newcommand{\Aut}{\mathrm{Aut}}
\newcommand{\PP}{\mathsf P}
\newcommand{\HH}{\mathsf H}
\newcommand{\cX}{\mathcal X}
\newcommand{\cC}{\mathcal C}
\newcommand{\cP}{\mathcal P}
\newcommand{\cS}{\mathcal S}
\newcommand{\cK}{\mathcal K}
\newcommand{\cJ}{\mathcal J}
\newcommand{\LC}{\mathbf{LC}}
\newcommand{\id}{\mathrm{id}}

\begin{document}

\title[G\"odel--Dummett logic over finite categories of models]{G\"odel--Dummett
logic over a finite category of models: the models of the translated
theory are strings of arrows, and a subgroup obstruction to commuting with the imaginaries}
\author{Ryota Shibusawa}
\thanks{Corresponding author. E-mail: \texttt{r-shibusawa@daiichi-koudai.ac.jp}}
\address{Daiichi Institute of Technology}
\email{r-shibusawa@daiichi-koudai.ac.jp}
\keywords{intermediate logic, G\"odel--Dummett logic, Heyting category, pretopos completion, imaginaries, Kripke model, EI-category, presheaf}
\subjclass[2020]{03G30 (primary), 03B55, 18F10, 18B25 (secondary)}

\begin{abstract}
Let $\cX$ be a finite category in which every endomorphism is
invertible, and $\cP$ the pretopos of finite set-valued functors on
$\cX$, a first-order theory whose models are the objects of $\cX$. For
an intermediate logic $L$ proving the linearity axiom $(p\to q)\vee
(q\to p)$ we describe the theory read in first-order logic over $L$:
its models are the strings of non-invertible arrows of $\cX$, of length
bounded by the depth of $L$, and its pretopos completion is the
category of finite set-valued functors on the category $\cS$ of
strings, whose arrows are the arrows of $\cX$ induced by isomorphisms
of final segments of the domain onto the codomain. The translation
into $L$ commutes with the imaginaries if and only if every finite
functor on $\cS$ embeds in one pulled back from $\cX$, a decidable
condition. It holds for classical logic and the logic of the
two-element chain. For the other logics it holds for finite posets,
for products of a finite poset and a finite group, and when $\cX$ has
only identity automorphisms; it fails when $\langle H\cap S_1,H\cap
S_2\rangle\ne H\cap\langle S_1,S_2\rangle$ for the group $H$ of the
automorphisms of the first object of a string which extend to the
string and the stabilisers $S_1,S_2$ of two arrows of $\cS$ into it.
In a family of categories with three levels it holds exactly when a
subgroup $A$ of an abelian group distributes over subgroups $T_j$,
$A\cap\sum_JT_j=\sum_J(A\cap T_j)$. The main tool is the linearity of
the canonical Kripke model with elements: two homomorphisms out of a
model of the translated theory are comparable.
\end{abstract}

\maketitle

\section{Introduction}

For an intermediate logic $L$ and a small coherent category $\cC$ let
$\HH_L\cC$ be the syntactic category of the internal coherent theory of
$\cC$ in first-order logic over $L$, and $\PP$ the pretopos
completion, which adds the imaginaries. There is a full and faithful
comparison
\[
c^L_\cC\colon\HH_L\PP\cC\longrightarrow\PP\HH_L\cC
\]
\cite{Interm}. For intuitionistic logic it is an equivalence only if
$\cC$ is Boolean \cite{Heyting}; for the logics proving the linearity
axiom it is an equivalence for all propositional theories
\cite{Interm}, and for the first-order theory of pure equality it is an
equivalence only for classical logic \cite{Interm}. This paper studies
the logics $L$ proving linearity over the first-order theories which
are closest to the propositional ones: those with finitely many models
up to isomorphism, all sorts being finite in each of them.

Let $\cX$ be a finite category and $\cP=\Fin^{\cX}$ the pretopos of
functors from $\cX$ to finite sets; if idempotents split in $\cX$, the
models of $\cP$ are the objects of $\cX$ and the homomorphisms are the
arrows. For a finite poset $\cX$ this is a finite propositional theory.
We assume that $\cX$ is an \emph{EI-category}: every endomorphism is an
isomorphism. A \emph{string} is a sequence $s=(m_0\xrightarrow{a_1}m_1
\to\dots\xrightarrow{a_r}m_r)$ of non-invertible arrows, $r\ge0$, and
$\cS_n$ is the category of the strings with at most $n$ objects, one
in each isomorphism class, an
arrow $s\to s'$ being an arrow $m_0\to m'_0$ of $\cX$ of the form
$\iota_0a_j\cdots a_1$ for an isomorphism of strings $\iota$ from the
final segment $(m_j\to\dots\to m_r)$ onto $s'$. Let $\pi\colon\cS_n\to
\cX$ be the first object.

\begin{thm}\label{thm:main}
Let $\cX$ be a finite EI-category, $\cP=\Fin^{\cX}$, and $L$ an
intermediate logic proving the linearity axiom, with $n=n_L\in
\{1,2,\dots,\infty\}$ the supremum of the lengths of the finite chains
which are frames of $L$. Then
\[
\PP\HH_L\cP\;\simeq\;\Fin^{\cS_n},
\]
the unit of $\HH_L$ followed by that of $\PP$ corresponding to
$\pi^*\colon\Fin^{\cX}\to\Fin^{\cS_n}$, and $\HH_L\cP$ to the full
subcategory of the subfunctors of the functors $\pi^*A$. Consequently
$c^L_\cP$ is an equivalence if and only if every functor $W\colon\cS_n
\to\Fin$ is a subfunctor of $\pi^*Y$ for some $Y\colon\cX\to\Fin$, if
and only if the unit $W\to\pi^*\mathrm{Lan}_\pi W$ of the left Kan
extension is injective for every $W$.
\end{thm}

For $n=1$, classical logic, $\cS_1$ is a skeleton of the groupoid of
the isomorphisms of $\cX$. The condition of the theorem is combinatorial.
We prove:
\begin{itemize}
\item it holds if for all arrows $v\colon t'\to t$, $u\colon t'\to s$
of $\cS_n$ every arrow $g$ of $\cX$ with $gv=u$ is an arrow $t\to s$ of
$\cS_n$ (Proposition~\ref{prop:T}); this is the case for finite posets,
which gives again the theorem of \cite{Interm} for finite propositional
theories, for the product of a finite poset and a finite group
(Corollary~\ref{cor:product}), and for all finite EI-categories if
$n_L\le2$ (Corollary~\ref{cor:two});
\item more generally it holds if, for every string $s$, the arrows of
$\cS_n$ into $s$ are a natural retract of the arrows of $\cX$ which are
connected to them (Proposition~\ref{prop:retract}); this is the case
if $\cX$ has no automorphisms other than identities
(Corollary~\ref{cor:rigid});
\item it fails if for a string $t$, with $H\le G=\Aut_{\cX}(m_0)$ the
group of the first components of its automorphisms, and two arrows of $\cS_n$ into $t$ with
stabilisers $S_1,S_2\le G$ one has $\langle H\cap S_1,H\cap S_2\rangle
\ne H\cap\langle S_1,S_2\rangle$ (Proposition~\ref{prop:modular}); this
happens in a category with four objects and fifteen arrows
(Example~\ref{ex:modular}), for every $L$ with $n_L\ge3$;
\item it fails if an arrow $g$ of $\cX$ which is not an arrow of
$\cS_n$ becomes one after composition with two arrows $v_1$, $v_2$ of
$\cS_n$, in two ways which are not connected inside $\cS_n$
(Proposition~\ref{prop:onestep}); a category with seventeen arrows
shows that this is not a consequence of the subgroup obstruction
(Example~\ref{ex:top});
\item for the categories with objects $0_1,\dots,0_m,1,2$ given by
subgroups $A,T_1,\dots,T_m$ of a finite abelian group it holds, for
$n_L\ge3$, if and only if $A\cap\sum_{j\in J}T_j=\sum_{j\in J}(A\cap
T_j)$ for every set $J$ of indices (Proposition~\ref{prop:dist}); among
them are a category with $28$ arrows for which the comparison is an
equivalence although there is no retraction, and, for every $m\ge3$, a
category for which it is not an equivalence although it is one for
every full subcategory omitting one of the objects $0_j$, and the
obstruction of Proposition~\ref{prop:onestep} is absent
(Example~\ref{ex:dist});
\item in general the condition has to be tested on finitely many
functors only, one for each string, so that it is decidable
(Proposition~\ref{prop:test}).
\end{itemize}
For classical logic the assumption on the endomorphisms is not needed:
the comparison is an equivalence for every finite category in which
idempotents split (Proposition~\ref{prop:classical}). The results are
summarised as follows, $\cX$ being a finite category in which
idempotents split.
\begin{center}
\begin{tabular}{lll}
logic $L$ & $\cX$ & $c^L_\cP$ is an equivalence\\\hline
classical ($n_L=1$) & any & always\\
linear, $n_L=2$ & EI & always\\
linear, $3\le n_L\le\infty$ & EI & iff the units $W\to\pi^*
\mathrm{Lan}_\pi W$ are injective;\\
&& not always\\
intuitionistic & any & only for groupoids \cite{Heyting}
\end{tabular}
\end{center}
So, in contrast with the propositional case, the
G\"odel--Dummett translation does not commute with the imaginaries of
all these theories, the obstructions are of a group-theoretic nature,
and they appear exactly from the three-element chain on. For
intuitionistic logic the comparison is an equivalence only
if $\cX$ is a groupoid \cite[Thm.~1.1]{Heyting}; the product of the
poset $0<1$ and the group $\mathbb Z/2$, for which an object outside
the image is given in \cite[Ex.~7.16]{Heyting}, is among the
categories for which the comparison is an equivalence for all $L$
proving linearity.

\subsection*{What is new and what is known}
New, to the author's knowledge: Theorem~\ref{thm:main}; the linearity
of the canonical model with elements (Lemma~\ref{lem:linear}) and the
description of the models by strings
(Proposition~\ref{prop:strings}); the classical case
(Proposition~\ref{prop:classical}); Propositions~\ref{prop:T},
\ref{prop:retract}, \ref{prop:modular}, \ref{prop:onestep},
\ref{prop:dist} and~\ref{prop:test} with
their corollaries, and Examples~\ref{ex:parallel}, \ref{ex:modular},
\ref{ex:top} and~\ref{ex:dist}. Known and used:
the translation $\HH_L$, the comparison and the frames of the logics
proving linearity \cite{Interm}; the completeness theorem of coherent
logic and the pretopos completion \cite{Elephant}; Joyal's completeness
theorem for Kripke models in the form of a truth lemma for the
category of coherent models, proved by the method of diagrams as in
\cite[Lemma~7.13]{Heyting}; and the fact that the prime filters above
a prime filter of a linear Heyting algebra form a chain \cite{Horn69},
of which Lemma~\ref{lem:linear} is the version with elements.
The description of the models by chains in the propositional case,
where $\cX$ is a finite poset, is known \cite{AGM08,Carai,Interm}.
EI-categories are standard in the representation theory of finite
categories; we use only the definition.

\section{The setting}\label{sec:setting}

Conventions are those of \cite{Interm,Heyting}. Let $\cX$ be a finite
category in which idempotents split, $\cP=\Fin^{\cX}$, $y(x)=\cX(x,-)$,
and $y(\gamma)\colon y(x')\to y(x)$ the composition with $\gamma\colon
x\to x'$. Every $A$ in $\cP$ is covered by the morphisms $q_a\colon
y(x)\to A$, $a\in A(x)$, $q_a(\varphi)=A(\varphi)(a)$. The models of
$\cP$ are, up to isomorphism, the evaluations at the objects of $\cX$,
and the homomorphisms between these are the arrows of $\cX$
\cite[\S7]{Heyting}.

Let $L$ be an intermediate logic and $\HH_L\cP$ the syntactic category
of the internal theory of $\cP$ in first-order logic over $L$
\cite[\S2]{Interm}: intuitionistic first-order logic with the
additional axioms $\lambda(\varphi_1,\dots,\varphi_k)$, for
$\lambda(p_1,\dots,p_k)\in L$ and first-order formulas $\varphi_i$ in
a common context, so that the formulas of $L$ are valid in every
lattice of subobjects of $\HH_L\cP$. Its objects are the subobjects
$U\le\eta A$, $A$ in $\cP$, and we omit $\eta$. The subobjects of the objects $A$ of $\cP$
in $\HH_L\cP$ are obtained from those in $\cP$ by finite meets and
joins, $\Rightarrow$, and $f^*$, $\exists_f$, $\forall_f$ for the
morphisms $f$ of $\cP$. Since $\cP$ is a pretopos, $\HH_L\PP\cP\simeq
\HH_L\cP$, and $c^L_\cP$ is the unit $\HH_L\cP\to\PP\HH_L\cP$; it is an
equivalence if and only if $\HH_L\cP$ is a pretopos
\cite[Lemma~2.2]{Interm}.

A \emph{model} of $\HH_L\cP$ is a coherent functor $\hat G\colon
\HH_L\cP\to\Set$; it is not required to be a Heyting functor. Its restriction to $\cP$ is isomorphic to the
evaluation at an object $m$; we fix such an isomorphism and regard
$\hat G(U)$, for $U\le A$, as a subset of $A(m)$, and say that $\hat G$
is \emph{over} $m$. A natural transformation $\hat G\to\hat G'$ of
models over $m$, $m'$ restricts to an arrow $\varphi\colon m\to m'$ and
is determined by it; so a \emph{homomorphism} $\hat G\to\hat G'$ is an
arrow $\varphi\colon m\to m'$ of $\cX$ with $A(\varphi)(\hat G(U))
\subseteq\hat G'(U)$ for all $U\le A$. Let $\cK$ be the category of the
models and homomorphisms, with one model in each isomorphism class; the
functor $\cK\to\cX$ is faithful.

\begin{lem}\label{lem:generic}
An arrow $\varphi\colon m\to m'$ is a homomorphism $\hat G\to\hat G'$
if and only if $\varphi\in\hat G'(U)$ for every $U\le y(m)$ with
$\id_m\in\hat G(U)$.
\end{lem}
\begin{proof}
If $b\in\hat G(U)$, $U\le A$, then $b=q_b(\id_m)$, so $\id_m\in\hat
G(q_b^*U)$, and $A(\varphi)(b)=q_b(\varphi)$; so $A(\varphi)(b)\in\hat
G'(U)$ iff $\varphi\in\hat G'(q_b^*U)$.
\end{proof}

\begin{lem}[truth lemma]\label{lem:truth}
Let $\hat G$ be a model over $m$, $f\colon B\to A$ a morphism of $\cP$,
$U,V\le B$ in $\HH_L\cP$, $\alpha\in A(m)$ and $\beta\in B(m)$.
\begin{enumerate}
\item[(a)] $\alpha\in\hat G(\forall_fU)$ if and only if for every
homomorphism $\varphi\colon\hat G\to\hat G'$ and every $\beta'\in
B(m')$ with $f(\beta')=A(\varphi)(\alpha)$ one has $\beta'\in\hat
G'(U)$.
\item[(b)] $\beta\in\hat G(U\Rightarrow V)$ if and only if for every
homomorphism $\varphi\colon\hat G\to\hat G'$, $B(\varphi)(\beta)\in
\hat G'(U)$ implies $B(\varphi)(\beta)\in\hat G'(V)$.
\end{enumerate}
\end{lem}
\begin{proof}
The implications from left to right hold because $\varphi$ preserves
$\forall_fU$ and $U\Rightarrow V$, and $f^*\forall_fU\le U$,
$(U\Rightarrow V)\wedge U\le V$. For the converse of (a) let
$\alpha\notin\hat G(\forall_fU)$, and consider the classical theory of
the coherent functors on $\HH_L\cP$ with a constant $c_b$ for every
element $b$ of the sets $\hat G(X)$, $X$ an object of $\HH_L\cP$, the
axioms $g(c_b)=c_{g(b)}$ for the morphisms $g$ of $\HH_L\cP$, a new
constant $d$ of sort $B$ and the axioms $f(d)=c_\alpha$ and $\neg
U(d)$. If it had no model, a finite part, written as a subobject
$\theta$ of a product of $A$ with finitely many objects, with a point
in $\hat G(\theta)$ over $\alpha$, would give $V_0:=\exists_{\pi_0}
\theta\le A$ with $\alpha\in\hat G(V_0)$ and $f^*V_0\le U$ by the
completeness theorem for the coherent category $\HH_L\cP$; then
$V_0\le\forall_fU$, contradicting $\alpha\notin\hat G(\forall_fU)$. A
model of the theory is a coherent functor $\hat G'$ with a natural
transformation $\hat G\to\hat G'$ and an element $\beta'\in B(m')
\setminus\hat G'(U)$ over the image of $\alpha$. This is the argument
of \cite[Lemma~7.13]{Heyting}. (b) is the same, with the axioms
$U(c_\beta)$ and $\neg V(c_\beta)$.
\end{proof}

So the evaluation $U\mapsto(\hat G\mapsto\hat G(U))$ is a Heyting
functor $\HH_L\cP\to\Set^{\cK}$, which is conservative by the
completeness theorem of coherent logic: if $\hat G(U)\subseteq\hat
G(V)$ for all models, then $U\le V$.

\subsection*{Kripke models}
Let $K$ be a finite poset such that every $\uparrow\!k$ is a frame of
$L$, and $M\colon K\to\cX$ a functor. Then $\Set^K$ is an $L$-category,
and $M$ extends to a Heyting functor $\HH_L\cP\to\Set^K$
\cite[Lemma~2.3]{Interm}; its value at $U\le A$ is the family of the
subsets of $A(M_k)$ defined by the forcing relation of $M$. The
evaluation at a node $k$ is a model of $\HH_L\cP$ over $M_k$, and the
transitions $M(k\le l)$ are homomorphisms between these models.

If $L$ proves the linearity axiom, the finite frames of $L$ are the
finite posets in which every $\uparrow\!k$ is a chain with at most
$n_L$ elements \cite[Lemma~3.1]{Interm}.

\section{Two homomorphisms out of a model are comparable}

\begin{lem}[linearity]\label{lem:linear}
Let $\cX$ be a finite category in which idempotents split, $L$ an
intermediate logic proving $(p\to q)\vee(q\to p)$, and $\varphi\colon
\hat G\to\hat G'$, $\psi\colon\hat G\to\hat G''$ homomorphisms of
models of $\HH_L\cP$. Then there is a homomorphism $\chi\colon\hat G'
\to\hat G''$ with $\chi\varphi=\psi$, or a homomorphism $\chi'\colon
\hat G''\to\hat G'$ with $\chi'\psi=\varphi$.
\end{lem}
\begin{proof}
Let the models be over $m$, $m'$, $m''$, and suppose that there is no
such $\chi$. The set of the arrows $\chi\colon m'\to m''$ of $\cX$ with
$\chi\varphi=\psi$ is finite, and none of them is a homomorphism; by
Lemma~\ref{lem:generic} there is for each of them $U_\chi\le y(m')$
with $\id_{m'}\in\hat G'(U_\chi)$ and $\chi\notin\hat G''(U_\chi)$. Let
$U'$ be the meet of the $U_\chi$ (the top element, if there is no such
arrow), and $E':=\exists_{y(\varphi)}U'\le y(m)$. Then $\varphi\in\hat
G'(E')$, since $\varphi=y(\varphi)(\id_{m'})$ and $\id_{m'}\in\hat
G'(U')$; and $\psi\notin\hat G''(E')$, since an element of $\hat
G''(U')$ over $\psi$ is an arrow $\chi$ with $\chi\varphi=\psi$ and
$\chi\in\hat G''(U')$. If there is no $\chi'$ either, we get in the same
way $E''\le y(m)$ with $\psi\in\hat G''(E'')$ and $\varphi\notin\hat
G'(E'')$. In $\HH_L\cP$ the subobject $(E'\Rightarrow E'')\vee
(E''\Rightarrow E')$ of $y(m)$ is the top element, so $\id_m$ lies in
$\hat G(E'\Rightarrow E'')$ or in $\hat G(E''\Rightarrow E')$. In the
first case $\varphi\in\hat G'(E'\Rightarrow E'')$, the homomorphism
$\varphi$ sending $\id_m$ to $\varphi$; with $\varphi\in\hat G'(E')$
this gives $\varphi\in\hat G'(E'')$, a contradiction. The second case
is symmetric.
\end{proof}

By Lemma~\ref{lem:truth} the forcing at a model $\hat G$, in the Kripke
model $\cK$, involves the homomorphisms out of $\hat G$;
Lemma~\ref{lem:linear} says that these are totally preordered by
$\varphi\preceq\psi$ iff $\psi=\chi\varphi$ for a homomorphism $\chi$.

\section{EI-categories: the models are strings}\label{sec:strings}

From now on $\cX$ is a finite EI-category and $L$ proves linearity. We
use the following consequences of the definition. A split monomorphism
or epimorphism of $\cX$ is invertible: if $\rho\gamma=\id$, then
$\gamma\rho$ is an idempotent endomorphism, hence invertible, hence the
identity. So idempotents split, and a composite $\delta\gamma$ is
invertible only if $\gamma$ and $\delta$ are. If there is a
non-invertible arrow $x\to x'$, then $x$ and $x'$ are not isomorphic;
so the isomorphism classes of objects are strictly partially ordered
by the existence of a non-invertible arrow. The non-invertible arrows
out of $x$ form a subfunctor $N_x\le y(x)$.

\begin{lem}\label{lem:iso}
A homomorphism of models of $\HH_L\cP$ whose underlying arrow is
invertible in $\cX$ is an isomorphism of models.
\end{lem}
\begin{proof}
We show, for all homomorphisms $\chi\colon\hat G'\to\hat G''$ with
invertible underlying arrow at once and all $U\le A$, that $b\in\hat
G'(U)$ if $A(\chi)(b)\in\hat G''(U)$, by induction on the construction
of $U$. This is clear for the subobjects in $\cP$ and compatible with
finite meets and joins and $f^*$; for $\exists_f$ it holds because
$\chi$ is bijective on every sort. Let $A(\chi)(b)\in\hat
G''(\forall_fU)$, $f\colon B\to A$; by Lemma~\ref{lem:truth} we have to
show $\beta\in\hat G'''(U)$ for a homomorphism $\lambda\colon\hat G'\to
\hat G'''$ and $\beta$ over $A(\lambda)(b)$. By Lemma~\ref{lem:linear}
there is a homomorphism $\mu$ with $\mu\chi=\lambda$ or with
$\mu\lambda=\chi$. In the first case $A(\lambda)(b)=A(\mu)A(\chi)(b)
\in\hat G'''(\forall_fU)$ and $\beta\in\hat G'''(U)$. In the second
case $\lambda$ is a split monomorphism, hence invertible, and so is
$\mu\colon\hat G'''\to\hat G''$; $B(\mu)(\beta)$ lies over
$A(\chi)(b)$, so $B(\mu)(\beta)\in\hat G''(U)$, and $\beta\in\hat
G'''(U)$ by the induction hypothesis for $\mu$ and $U$. The case of
$U\Rightarrow V$ is the same.
\end{proof}

A \emph{string} is a sequence $s=(m_0\xrightarrow{a_1}m_1\to\dots
\xrightarrow{a_r}m_r)$ of non-invertible arrows of $\cX$, $r\ge0$; its
\emph{length} is $r+1$, its \emph{final segments} are the strings
$s_{\ge j}=(m_j\to\dots\to m_r)$, and an isomorphism of strings is a
family of isomorphisms of the objects commuting with the arrows. The
length of a string is at most the number of isomorphism classes of
objects of $\cX$. The string $s$ is a Kripke model over the chain
$0<\dots<r$, with $M_j=m_j$; if $r+1\le n_L$ let $\hat G_s$ be the
model of $\HH_L\cP$ given by the evaluation at the node $0$. The
evaluation at the node $j$ is $\hat G_{s_{\ge j}}$, and
$a_j\cdots a_1\colon\hat G_s\to\hat G_{s_{\ge j}}$ is a homomorphism.

\begin{prop}\label{prop:strings}
Let $\cX$ be a finite EI-category, $L$ an intermediate logic proving
linearity and $n=n_L$.
\begin{enumerate}
\item[(a)] Every model of $\HH_L\cP$ is isomorphic to $\hat G_s$ for a
string $s$ of length at most $n$.
\item[(b)] The homomorphisms $\hat G_s\to\hat G_{s'}$ are the arrows
$\iota_0a_j\cdots a_1$, for $0\le j\le r$ and $\iota$ an isomorphism of
strings from $s_{\ge j}$ onto $s'$; the index $j$ is determined by
$s'$. In particular $\hat G_s\cong\hat G_{s'}$ if and only if $s$ and
$s'$ are isomorphic strings.
\end{enumerate}
\end{prop}
\begin{proof}
\emph{Step 1: the string of a model.} Let $\hat G$ be a model over
$m_0$. By Lemma~\ref{lem:linear} the homomorphisms out of $\hat G$ are
totally preordered by $\preceq$. If $\varphi\preceq\psi\preceq
\varphi$, say $\psi=\chi\varphi$ and $\varphi=\chi'\psi$ with
$\chi\colon\hat G'\to\hat G''$, then $\chi'\chi$ is an endomorphism of
$m'$ and $\chi\chi'$ one of $m''$; so $\chi$ is invertible in $\cX$,
and an isomorphism of models by Lemma~\ref{lem:iso}. If $\varphi\prec
\psi=\chi\varphi$, then $\chi$ is not an isomorphism of models, hence
not invertible in $\cX$, and the isomorphism class of the codomain of
$\psi$ is strictly above that of $\varphi$. So there are finitely many
classes; choose representatives $\id=\varphi_0\prec\varphi_1\prec\dots
\prec\varphi_r$, $\varphi_j\colon\hat G\to\hat G_j$, and homomorphisms
$a_j\colon\hat G_{j-1}\to\hat G_j$ with $\varphi_j=a_j\varphi_{j-1}$.
The $a_j$ are not invertible, and $s:=(m_0\xrightarrow{a_1}\dots
\xrightarrow{a_r}m_r)$ is a string.

\emph{Step 2: every homomorphism $\lambda$ out of $\hat G_j$ is
$\chi a_{j'}\cdots a_{j+1}$ for some $j'\ge j$ and an isomorphism of
models $\chi$.} The homomorphism $\lambda\varphi_j$ out of $\hat G$ is
equivalent to some $\varphi_{j'}$, and $j'\ge j$ because $\varphi_j
\preceq\lambda\varphi_j$; so $\lambda\varphi_j=\iota a_{j'}\cdots
a_{j+1}\varphi_j$ with an isomorphism of models $\iota$. By
Lemma~\ref{lem:linear} at $\hat G_j$, $\lambda$ and $\iota a_{j'}
\cdots a_{j+1}$ differ by a homomorphism between their codomains in
one of the two directions; it is an endomorphism in $\cX$, hence
invertible, hence an isomorphism of models.

\emph{Step 3: $\hat G_j(U)$ is the value of $U$ at the node $j$ of the
Kripke model $s$, for all $j$ and all $U\le A$,} where the forcing
relation of $s$ is defined by the usual clauses, whether or not the
chain is a frame of $L$. By induction on the construction of $U$: both
sides agree on the subobjects in $\cP$ and are compatible with the
coherent operations; for $\forall_fU$, by Lemma~\ref{lem:truth} and
Step 2, $\alpha\in\hat G_j(\forall_fU)$ iff for all $j'\ge j$, all
isomorphisms of models $\chi$ out of $\hat G_{j'}$ and all $\beta$
over $A(\chi a_{j'}\cdots a_{j+1})(\alpha)$ we have $\beta$ in the
value of $U$ at the codomain of $\chi$, iff, applying $\chi^{-1}$, for
all $j'\ge j$ and all $\beta$ over $A(a_{j'}\cdots a_{j+1})(\alpha)$
we have $\beta\in\hat G_{j'}(U)$, which by induction is the forcing
clause at the node $j$. The case of $\Rightarrow$ is the same.

\emph{Step 4: $r+1\le n$.} Let $U_i\le y(m_0)$ be the image of
$y(a_i\cdots a_1)$, for $1\le i\le r$, the arrows out of $m_0$ which
factor through $a_i\cdots a_1$. The arrow $a_j\cdots a_1$ lies in $U_i$
iff $j\ge i$: if $a_j\cdots a_1=\beta a_i\cdots a_1$ with $j<i$, then
$\beta a_i\cdots a_{j+1}$ is an endomorphism of $m_j$, so $a_{j+1}$
would be a split monomorphism. So at the element $\id_{m_0}$ the $U_i$
are forced exactly at the nodes $j\ge i$, and every valuation of
propositional variables by up-sets of the chain $0<\dots<r$ is
obtained from the $U_i$, the top and the bottom. For a formula
$\lambda(p_1,\dots,p_k)$ of $L$ and such a valuation, the subobject
$\lambda(U_{i_1},\dots,U_{i_k})$ of $y(m_0)$ is the top element in
$\HH_L\cP$, so $\id_{m_0}$ lies in its value under $\hat G$; by Step 3
this value is computed by forcing in the chain, where at the element
$\id_{m_0}$ it is the propositional forcing of $\lambda$. So $\lambda$
is valid in the chain, the chain is a frame of $L$, and $r+1\le n$.
Hence $\hat G_s$ is defined, and $\hat G\cong\hat G_s$ by Step 3; this
is (a).

\emph{Step 5: (b).} Let $s$ be a string of length at most $n$. We show
first that every homomorphism $\psi\colon\hat G_s\to\hat G'$ is
$\chi a_i\cdots a_1$ with $\chi$ an isomorphism of models. Let $i$ be
maximal with $a_i\cdots a_1\preceq\psi$, and $\psi=\lambda a_i\cdots
a_1$ with $\lambda\colon\hat G_{s_{\ge i}}\to\hat G'$; it suffices to
show that $\lambda$ is invertible in $\cX$. If $i=r$, the node $r$ is
the top of the chain, so the forcing there is classical and $\id_{m_r}
\in\hat G_{s_{\ge r}}(\neg N_{m_r})$; then $\lambda\in\hat G'(\neg
N_{m_r})$, so $\lambda\notin N_{m_r}(m')$. Let $i<r$ and suppose that
$\lambda$ is not invertible. By Lemma~\ref{lem:linear} and the
maximality of $i$ there is a homomorphism $\mu$ with $\mu\lambda=
a_{i+1}$. At the node $i$ of $s$ the element $\id_{m_i}$ forces
$N_{m_i}\Rightarrow I$, where $I\le y(m_i)$ is the image of
$y(a_{i+1})$: at the nodes $j>i$ it has become $a_j\cdots a_{i+1}$.
So $\lambda\in\hat G'(N_{m_i}\Rightarrow I)$, and as $\lambda\in
N_{m_i}$, $\lambda=\beta a_{i+1}$ for an arrow $\beta$ of $\cX$. Then
$a_{i+1}=\mu\beta a_{i+1}$, the endomorphism $\mu\beta$ is invertible,
so $\beta$ and $\mu$ are invertible, $\mu$ is an isomorphism of models
and $a_{i+1}\preceq\lambda$, contradicting the maximality of $i$.

It remains to show that the isomorphisms of models $\chi_0\colon\hat
G_u\to\hat G_t$ are the first components of the isomorphisms of
strings $u\to t$. An isomorphism of strings is an isomorphism of
Kripke models and preserves forcing. Conversely let $\chi_0$ be an
isomorphism of models, $u=(m_0\xrightarrow{a_1}\cdots)$ and
$t=(m'_0\xrightarrow{b_1}\cdots)$. It induces a bijection between the
homomorphisms out of $\hat G_u$ and out of $\hat G_t$ which preserves
$\preceq$; by what was shown, the classes are those of the $a_i\cdots
a_1$ and of the $b_i\cdots b_1$, which are pairwise inequivalent since
their codomains are pairwise non-isomorphic. So $u$ and $t$ have the
same length, and $b_1\chi_0=\chi_1a_1$ for an isomorphism of models
$\chi_1\colon\hat G_{u_{\ge1}}\to\hat G_{t_{\ge1}}$; by induction on
the length $\chi_1$ extends to an isomorphism of strings. The index
$j$ in (b) is determined because the codomains of the $a_j\cdots a_1$
are pairwise non-isomorphic.
\end{proof}

\begin{defi}
$\cS_n=\cS_n(\cX)$ is the category whose objects are the strings of
length at most $n$, one in each isomorphism class, and whose arrows
$s\to s'$ are the arrows $\iota_0a_j\cdots a_1\colon m_0\to m'_0$ of
$\cX$ of Proposition~\ref{prop:strings}(b); $\pi\colon\cS_n\to\cX$ is
the first object, a faithful functor. We write $\Aut_{\cS}(s)\le
\Aut_{\cX}(m_0)$ for the group of the first components of the
automorphisms of the string $s$.
\end{defi}

By Proposition~\ref{prop:strings} the category $\cK$ of models of
$\HH_L\cP$ is $\cS_n$, a finite category. For a poset $\cX$ the strings
are the strict chains and $\cS_n$ is the poset of the strict chains
with at most $n$ elements, ordered by final segments, of
\cite[\S6]{Interm}; this description of the propositional case by
chains is known, from the free G\"odel algebras over finite
distributive lattices \cite{AGM08} and over bounded distributive
lattices in the subvarieties \cite{Carai}.

\section{The pretopos completion of the translated theory}

\begin{lem}\label{lem:definable}
Let $A$ be in $\cP$ and $S$ a subfunctor of $\pi^*A\colon\cS_n\to\Fin$.
Then there is a unique $V\le A$ in $\HH_L\cP$ with $\hat G_s(V)=S(s)$
for all $s$.
\end{lem}
\begin{proof}
Uniqueness is the conservativity of the evaluation. For a string $s$
with first object $m$, a string $t$ and an arrow $\varphi\colon m\to
t_0$ which is not an arrow $s\to t$ of $\cS_n$ there is, by
Lemma~\ref{lem:generic}, $U\le y(m)$ with $\id_m\in\hat G_s(U)$ and
$\varphi\notin\hat G_t(U)$. Let $\tau_s\le y(m)$ be the meet of these
finitely many $U$. Then $\varphi\in\hat G_t(\tau_s)$ if and only if
$\varphi$ is an arrow $s\to t$ of $\cS_n$. Let $V$ be the join of the
$\exists_{q_a}\tau_s$ over the strings $s$ and the $a\in S(s)$. Then
$\hat G_t(\exists_{q_a}\tau_s)$ is the set of the $A(\varphi)(a)$,
$\varphi\colon s\to t$ in $\cS_n$, which is contained in $S(t)$ and
contains $a$ if $t=s$.
\end{proof}

\begin{lem}\label{lem:PP}
Let $F\colon\mathcal H\to\mathcal E$ be a coherent functor from a small
coherent category to a pretopos which is conservative and full on
subobjects, and such that every object of $\mathcal E$ is a quotient of
a finite coproduct of objects $F(X)$. Then the coherent functor
$\PP\mathcal H\to\mathcal E$ extending $F$ is an equivalence.
\end{lem}
\begin{proof}
An object of $\PP\mathcal H$ is a quotient $(\coprod_iX_i)/R$ with $R$
given by subobjects $R_{ij}\le X_i\times X_j$ in $\mathcal H$, and a
morphism to $(\coprod_jX'_j)/R'$ is a family of subobjects of the
$X_i\times X'_j$ which is a functional relation compatible with $R$
and $R'$; the same holds in $\mathcal E$ for the quotients of the
coproducts of the $F(X_i)$. Since $F$ is full on subobjects, every
such family in $\mathcal E$ comes from $\mathcal H$, and since $F$ is
conservative, the conditions on the families, which are inequalities
between subobjects built by the coherent operations, hold in
$\mathcal H$ if they hold in $\mathcal E$, and two families are equal
if their images are. So the extension is full and faithful, and an
equivalence relation on $\coprod_iF(X_i)$ in $\mathcal E$ comes from
one in $\PP\mathcal H$; hence the extension is essentially surjective.
\end{proof}

\begin{proof}[Proof of Theorem~\ref{thm:main}]
By Lemma~\ref{lem:truth} and Proposition~\ref{prop:strings} the
evaluation $F\colon\HH_L\cP\to\Fin^{\cS_n}$, $F(U)(s)=\hat G_s(U)$, is
a conservative Heyting functor with $F(\eta A)=\pi^*A$. It is full on
subobjects: a subfunctor of $F(U)$, $U\le A$, is a subfunctor of
$\pi^*A$, hence $F(V)$ for some $V\le A$ by Lemma~\ref{lem:definable},
and $V\le U$ by conservativity. Every $W\colon\cS_n\to\Fin$ is a
quotient of a finite coproduct of representable functors $\cS_n(s,-)$,
and $\cS_n(s,-)$ is a subfunctor of $\pi^*y(m_0)$, $\pi$ being
faithful. By Lemma~\ref{lem:PP}, $\PP\HH_L\cP\simeq\Fin^{\cS_n}$, and
$\HH_L\cP$ corresponds to the full subcategory of the subfunctors of
the $\pi^*A$. So $\HH_L\cP$ is a pretopos, that is $c^L_\cP$ is an
equivalence, if and only if every $W$ is isomorphic to a subfunctor of
some $\pi^*A$. A morphism $W\to\pi^*Y$ factors through the unit
$W\to\pi^*\mathrm{Lan}_\pi W$, and $\mathrm{Lan}_\pi W$ has finite
values, $\cS_n$ being finite; so $W$ embeds in some $\pi^*Y$ if and
only if the unit is injective.
\end{proof}

For classical logic the models can be described without the assumption
on the endomorphisms.

\begin{prop}[classical logic]\label{prop:classical}
Let $\cX$ be a finite category in which idempotents split, $\cX^\times$
the groupoid of its isomorphisms and $L$ classical logic. Then the
models of $\HH_L\cP$ are the objects of $\cX$, the homomorphisms are
the isomorphisms of $\cX$, $\PP\HH_L\cP\simeq\Fin^{\cX^\times}$, and
$c^L_\cP$ is an equivalence.
\end{prop}
\begin{proof}
The evaluation at an object $m$ is a coherent functor $\cP\to\Set$, and
$\Set$ is Boolean; so it extends to a Heyting functor on $\HH_L\cP$,
uniquely up to isomorphism \cite[\S2]{Interm}. Since $\HH_L\cP$ is
Boolean, every coherent functor on it is Heyting
\cite[Lemma~2.1]{Heyting}; so there is exactly one model over each $m$,
and the isomorphisms of $\cX$ are isomorphisms of models. Let
$\varphi\colon m\to m'$ be a homomorphism of models. It preserves
complements, hence reflects the subobjects of $\HH_L\cP$. For an object
$x$ of $\cX$ and $k\ge1$ let $E_k\le1$ be the image of the complement,
in $y(x)^k$, of the union of the diagonals; it is true in the model
over $m$ iff $\cX(x,m)$ has at least $k$ elements. As $\varphi$
preserves and reflects the $E_k$ and reflects the diagonal of $y(x)$,
the map $\cX(x,m)\to\cX(x,m')$ is an injection between sets of the
same cardinality, for every $x$; so $\varphi$ is invertible. Hence the
category $\cK$ of models is a skeleton of $\cX^\times$. The proofs of
Lemma~\ref{lem:definable} and of the equivalence $\PP\HH_L\cP\simeq
\Fin^{\cK}$ in Theorem~\ref{thm:main} used only
Lemma~\ref{lem:generic}, the conservativity of the evaluation, the
finiteness of $\cK$ and the faithfulness of $\pi\colon\cK\to\cX$, so
they apply. Finally let $W\colon\cK\to\Fin$
and $m$ in $\cK$; $\mathrm{Lan}_\pi W(m)$ is the colimit of $W(m')$ over
the category of the pairs $(m',g\colon m'\to m)$, whose arrows
$(m',g'v)\to(m'',g')$ are the isomorphisms $v$. In it $g'v$ is
invertible iff $g'$ is, so the pairs with $g$ invertible form a union
of connected components, with the terminal object $(m,\id)$, and the
unit $W(m)\to\mathrm{Lan}_\pi W(m)$ is injective.
\end{proof}

\section{When the translation commutes with the imaginaries}

For a string $s$ with first object $m$ let $\cJ_s$ be the category of
the pairs $(t,g)$, $t$ a string and $g\colon t_0\to m$ an arrow of
$\cX$, an arrow $(t,g)\to(t',g')$ being an arrow $v\colon t\to t'$ of
$\cS_n$ with $g'v=g$. Then $\mathrm{Lan}_\pi W(m)$ is the colimit of
$W(t)$ over $\cJ_s$, and the unit at $s$ is the map from $W(s)$, at the
object $(s,\id)$. Let $\cJ^0_s$ be the full subcategory of the $(t,g)$
with $g$ an arrow $t\to s$ of $\cS_n$; it has the terminal object
$(s,\id)$, so the colimit of $W$ over $\cJ^0_s$ is $W(s)$. An arrow of
$\cJ_s$ into an object of $\cJ^0_s$ has its domain in $\cJ^0_s$. Let
$\cJ^c_s\supseteq\cJ^0_s$ be the connected component of $(s,\id)$ in
$\cJ_s$, and $C_s(t)$ the set of the arrows $g$ with $(t,g)$ in
$\cJ^c_s$; so $\cS_n(t,s)\subseteq C_s(t)\subseteq\cX(t_0,s_0)$, and
$gv\in C_s(t')$ iff $g\in C_s(t)$, for $v\colon t'\to t$ in $\cS_n$.

\begin{prop}\label{prop:T}
Suppose that for all arrows $v\colon t'\to t$ and $u\colon t'\to s$ of
$\cS_n$, every arrow $g\colon t_0\to s_0$ of $\cX$ with $gv=u$ is an
arrow $t\to s$ of $\cS_n$. Then $c^L_\cP$ is an equivalence.
\end{prop}
\begin{proof}
The hypothesis says that an arrow of $\cJ_s$ out of an object of
$\cJ^0_s$ has its codomain in $\cJ^0_s$. So $\cJ^0_s$ is a union of
connected components of $\cJ_s$, the colimit of $W$ over $\cJ_s$ is the
disjoint union of $W(s)$ and the colimit over the other components,
and the unit is injective.
\end{proof}

\begin{cor}\label{cor:product}
Let $X$ be a finite poset and $\Gamma$ a finite group, regarded as a
category with one object. For $\cX=X\times\Gamma$, in particular for
$\cX=X$, the comparison $c^L_\cP$ is an equivalence for every $L$
proving linearity.
\end{cor}
\begin{proof}
An arrow of $\cX$ is a pair $(x\le x',\gamma)$; it is invertible iff
$x=x'$, so $\cX$ is an EI-category, and a string is a strict chain
$x_0<\dots<x_r$ with elements $\gamma_1,\dots,\gamma_r$. For strings
$t$, $s$ whose chains are $x_0<\dots<x_r$ and $x_j<\dots<x_r$, every
$\gamma_0\in\Gamma$ is the first component of a unique isomorphism of
strings from $t_{\ge j}$ onto $s$, the other components being
determined successively. So the arrows $t\to s$ of $\cS_n$ are all the
arrows $t_0\to s_0$ of $\cX$ if the chain of $s$ is a final segment of
that of $t$, and there are none otherwise. Let $v\colon t'\to t$,
$u\colon t'\to s$ and $g\colon t_0\to s_0$ with $gv=u$. The chains of
$t$ and $s$ are final segments of that of $t'$, and the first element
of $t$ is below that of $s$; so the chain of $s$ is a final segment of
that of $t$, and $g$ is an arrow of $\cS_n$. Proposition~\ref{prop:T}
applies.
\end{proof}

For $\Gamma=1$ Corollary~\ref{cor:product} is the theorem of \cite{Interm} that the translation
into a logic proving linearity commutes with the imaginaries of the
finite propositional theories. For $X=\{0<1\}$ and $\Gamma=\mathbb Z/2$
the pretopos $\cP$ is that of \cite[Ex.~7.16]{Heyting}, a proposition
$p$ and a sort with exactly two elements and decidable equality; over
intuitionistic logic a twisted form of this sort is not in the image
of the comparison, and it needs a Kripke model over a fork.

\begin{cor}\label{cor:two}
If $n_L\le2$, in particular for classical logic and for the logic of
the two-element chain, $c^L_\cP$ is an equivalence for every finite
EI-category $\cX$.
\end{cor}
\begin{proof}
Let $v\colon t'\to t$ and $u\colon t'\to s$ be arrows of $\cS_n$ and
$gv=u$. If $v$ is invertible, $g=uv^{-1}$ is an arrow of $\cS_n$.
Otherwise $t'=(m_0\xrightarrow{a}m_1)$, $t$ has length one and
$v=\iota a$ with $\iota\colon m_1\to t_0$ invertible. If $u$ were
invertible, $a$ would be a split monomorphism; so $u=\iota'a$ with
$\iota'\colon m_1\to s_0$ invertible, and $s$ has length one. Then
$\iota'^{-1}g\iota$ is an endomorphism of $m_1$, so $g$ is invertible,
and an invertible arrow between strings of length one is an arrow of
$\cS_n$. Proposition~\ref{prop:T} applies.
\end{proof}

The hypothesis of Proposition~\ref{prop:T} says that $C_s(t)=\cS_n(t,
s)$. It is the simplest case of the following.

\begin{prop}[retraction]\label{prop:retract}
Suppose that for every string $s$ of length at most $n$ there are maps
$r_t\colon C_s(t)\to\cS_n(t,s)$, $t$ in $\cS_n$, with $r_t(\kappa)=
\kappa$ for $\kappa\in\cS_n(t,s)$ and $r_{t'}(gv)=r_t(g)v$ for $g\in
C_s(t)$ and $v\colon t'\to t$ in $\cS_n$. Then $c^L_\cP$ is an
equivalence.
\end{prop}
\begin{proof}
The colimit of $W$ over $\cJ_s$ is the disjoint union of the colimits
over the connected components. On the colimit over $\cJ^c_s$ let
$\rho$ send the class of $w\in W(t)$ at the object $(t,g)$ to
$W(r_t(g))(w)\in W(s)$. This is well defined: for $v\colon(t',gv)\to
(t,g)$ and $w'\in W(t')$ the elements $w'$ at $(t',gv)$ and $W(v)w'$
at $(t,g)$ are both sent to $W(r_t(g)v)(w')$. Since $r_s(\id)=\id$,
$\rho$ is a left inverse of the map from $W(s)$ to the colimit over
$\cJ^c_s$, which is a summand of $\mathrm{Lan}_\pi W(s_0)$; so the unit
is injective.
\end{proof}

\begin{cor}\label{cor:rigid}
If every automorphism of $\cX$ is an identity, then $c^L_\cP$ is an
equivalence for every $L$ proving linearity.
\end{cor}
\begin{proof}
Two arrows $t\to s$ of $\cS_n$ are $\iota_0p$ and $\iota'_0p$ for the
same composite $p$ of arrows of $t$ and isomorphisms of strings
$\iota$, $\iota'$ onto $s$; these differ by an automorphism of the
string $s$, whose first component is the identity by assumption. So $\cS_n(t,s)$ has
at most one element. We show that it is non-empty if $C_s(t)$ is;
then the unique maps $r_t\colon C_s(t)\to\cS_n(t,s)$ satisfy the
conditions of Proposition~\ref{prop:retract}. It suffices to show, for
$v\colon t'\to t$ in $\cS_n$ and $g\colon t_0\to s_0$ in $\cX$ with
$\cS_n(t',s)\ne\emptyset$, that $\cS_n(t,s)\ne\emptyset$, the converse
being clear. Up to isomorphism $t=t'_{\ge j}$ and $s=t'_{\ge k}$, with
$t'=(m_0\to\dots\to m_r)$. If $k<j$, the composite $p\colon m_k\to m_j$
of the arrows of $t'$ is not invertible, and $gp$ is an endomorphism
of $m_k$, so $p$ would be a split monomorphism. So $k\ge j$, and $s$
is a final segment of $t$.
\end{proof}

\begin{ex}\label{ex:parallel}
Let $\cX$ have the objects $0$, $1$, $2$ and, besides the identities,
the arrows $a\colon0\to1$, $c,c'\colon1\to2$ and $d=ca=c'a\colon0\to2$:
\[
\begin{tikzcd}
0\arrow[r,"a"]&1\arrow[r,shift left,"c"]\arrow[r,shift right,"c'"']&2
\end{tikzcd}
\]
By Corollary~\ref{cor:rigid}, $c^L_\cP$ is an equivalence for every $L$
proving linearity. The hypothesis of Proposition~\ref{prop:T} fails for
$n\ge3$: for $v=a\colon(0\xrightarrow{a}1\xrightarrow{c}2)\to
(1\xrightarrow{c}2)$ and $u=d\colon(0\xrightarrow{a}1\xrightarrow{c}2)
\to(2)$ one has $c'v=u$, and $c'$ is not an arrow
$(1\xrightarrow{c}2)\to(2)$ of $\cS_n$. The retraction sends $c'$ to
$c$.
\end{ex}

\begin{prop}[subgroup obstruction]\label{prop:modular}
Let $t$ be a string of length at most $n$ with first object $m$,
$G=\Aut_{\cX}(m)$, $H=\Aut_{\cS}(t)\le G$, and let $v_i\colon t'_i\to
t$, $i=1,2$, be arrows of $\cS_n$ with stabilisers $S_i=\{g\in G\mid
gv_i=v_i\}$. If
\[
\langle H\cap S_1,\ H\cap S_2\rangle\ \ne\ H\cap\langle S_1,S_2\rangle,
\]
then $c^L_\cP$ is not an equivalence.
\end{prop}
\begin{proof}
Let $W$ be the pushout in $\Fin^{\cS_n}$ of the morphisms
$\cS_n(t,-)\to\cS_n(t'_i,-)$, $\kappa\mapsto\kappa v_i$, and $x_i\in
W(t'_i)$ the class of the identity; so $W(v_1)x_1=W(v_2)x_2=:y$. The
arrows $t'_i\to t$ of $\cS_n$ are the $hv_i$, $h\in H$: by
Proposition~\ref{prop:strings}(b) the index of the final segment is
determined, and two isomorphisms of strings onto $t$ differ by an
automorphism of $t$. So $W(t)$ is the quotient of $Hv_1\sqcup Hv_2$ by
the relation generated by $hv_1\sim hv_2$, and $hv_i=h'v_i$ iff
$h^{-1}h'\in S_i$; hence the stabiliser of $y$ in $H$, for the action
of $H$ on $W(t)$ by the arrows $t\to t$, is $\langle H\cap S_1,H\cap
S_2\rangle$. Suppose that $\theta\colon W\to\pi^*Y$ is injective, and
let $\bar y=\theta_t(y)\in Y(m)$. By naturality $\bar y=Y(v_i)
\theta(x_i)$, so $Y(g)\bar y=\bar y$ for $g\in S_i$, hence for
$g\in\langle S_1,S_2\rangle$. For $h\in H\cap\langle S_1,S_2\rangle$
we get $\theta_t(W(h)y)=Y(h)\bar y=\theta_t(y)$, so $W(h)y=y$ and
$h\in\langle H\cap S_1,H\cap S_2\rangle$. So under the hypothesis $W$
is not a subfunctor of any $\pi^*Y$, and Theorem~\ref{thm:main}
applies.
\end{proof}

\begin{ex}\label{ex:modular}
Let $G=\mathbb Z/2\times\mathbb Z/2$, $S_1=\langle(1,0)\rangle$,
$S_2=\langle(0,1)\rangle$ and $H=\langle(1,1)\rangle$. Let $\cX$ have
the objects $0$, $0'$, $1$, $2$, with $\Aut(1)=G$ and trivial
automorphism groups otherwise, and the arrows
\begin{gather*}
\cX(0,1)=G/S_1,\qquad\cX(0',1)=G/S_2,\qquad\cX(1,2)=H\backslash G,\\
\cX(0,2)=\{d\},\qquad\cX(0',2)=\{d'\},
\end{gather*}
composition being given by the actions of $G$ on the cosets and by the
unique maps to the one-element sets:
\[
\begin{tikzcd}[column sep=large]
0\arrow[dr,"G/S_1"]&&\\
&1\arrow[r,"H\backslash G"]\arrow[loop above,"G"]&2\\
0'\arrow[ur,"G/S_2"']&&
\end{tikzcd}
\]
This is a category with fifteen arrows, and an EI-category. Let
$a=S_1\in\cX(0,1)$, $a'=S_2\in\cX(0',1)$ and $c=H\in\cX(1,2)$. The
string $t=(1\xrightarrow{c}2)$ has $\Aut_{\cS}(t)=\{g\mid cg=c\}=H$,
and the arrows $a\colon(0\xrightarrow{a}1\xrightarrow{c}2)\to t$ and
$a'\colon(0'\xrightarrow{a'}1\xrightarrow{c}2)\to t$ of $\cS_3$ have
the stabilisers $S_1$ and $S_2$. Here $H\cap S_1=H\cap S_2=1$ and
$\langle S_1,S_2\rangle=G$; so by Proposition~\ref{prop:modular} the
comparison $c^L_\cP$ is not an equivalence for the logics $L$ proving
linearity with $n_L\ge3$, in particular for the G\"odel--Dummett
logic. In terms of Kripke models: the functor $W$ of the proof has two
points at a node of type $t$, exchanged by $(1,1)$; a point of this
kind which comes from a node below along $a$ and from a node below
along $a'$ must, in an imaginary of $\cP$, be fixed by $(1,0)$ and by
$(0,1)$.
\end{ex}

The same construction gives an obstruction which does not refer to
automorphisms of the first object.

\begin{prop}[one-step obstruction]\label{prop:onestep}
Let $s$, $t$ be strings of length at most $n$, $g\colon t_0\to s_0$ an
arrow of $\cX$, and $v_i\colon t_i\to t$, $i=1,2$, arrows of $\cS_n$
such that $gv_1$ and $gv_2$ are arrows $t_i\to s$ of $\cS_n$. Let
$\sim$ be the equivalence relation on $\cS_n(t_1,s)\sqcup\cS_n(t_2,s)$
generated by $\kappa v_1\sim\kappa v_2$, $\kappa\in\cS_n(t,s)$. If
$gv_1\not\sim gv_2$, then $c^L_\cP$ is not an equivalence.
\end{prop}
\begin{proof}
Let $W$ be the pushout of the morphisms $\cS_n(t,-)\to\cS_n(t_i,-)$,
$\kappa\mapsto\kappa v_i$, and $x_i\in W(t_i)$ the classes of the
identities, so that $W(v_1)x_1=W(v_2)x_2=:y$. Then $W(s)$ is the
quotient of $\cS_n(t_1,s)\sqcup\cS_n(t_2,s)$ by $\sim$, and the
elements $z_i:=W(gv_i)x_i$, the classes of the $gv_i$, are different.
For a morphism $\theta\colon W\to\pi^*Y$ we have
$\theta_s(z_i)=Y(gv_i)\theta_{t_i}(x_i)=Y(g)Y(v_i)\theta_{t_i}(x_i)=
Y(g)\theta_t(y)$ for $i=1,2$; so $\theta$ is not injective, and
Theorem~\ref{thm:main} applies.
\end{proof}

If $g$ is an arrow $t\to s$ of $\cS_n$, then $gv_1\sim gv_2$; so the
hypothesis implies that $(t,g)$ is an object of $\cJ^c_s$ outside
$\cJ^0_s$, and that there is no retraction as in
Proposition~\ref{prop:retract}, which would give $gv_i=r_t(g)v_i$.

\begin{ex}\label{ex:top}
Let $G=\mathbb Z/2\times\mathbb Z/2$, $a=(1,0)$, $b=(0,1)$. Let $\cX$
have the objects $0$, $0'$, $1$, $2$, with $\Aut(2)=\langle b\rangle$
and trivial automorphism groups otherwise, and the arrows
\begin{gather*}
\cX(0,1)=\cX(0,2)=G/\langle a\rangle,\qquad\cX(1,2)=G,\\
\cX(0',1)=\cX(0',2)=G/\langle a+b\rangle,
\end{gather*}
composition being the addition of $G$. This is an EI-category with
seventeen arrows. Let $s=(2)$, $t=(1\xrightarrow{0}2)$, so that
$\cS_3(t,s)=\{0,b\}$, and $g=a\in\cX(1,2)$. Let $v_1=\langle a\rangle
\in\cX(0,1)$ and $v_2=\langle a+b\rangle\in\cX(0',1)$, arrows of
$\cS_3$ into $t$ from the strings $t_1=(0\xrightarrow{v_1}1
\xrightarrow{0}2)$ and $t_2=(0'\xrightarrow{v_2}1\xrightarrow{0}2)$.
Then $gv_1=a+\langle a\rangle=0\cdot v_1$ and $gv_2=a+\langle a+b
\rangle=b\cdot v_2$ are arrows $t_i\to s$ of $\cS_3$. The relation
$\sim$ is generated by $0\cdot v_1\sim0\cdot v_2$ and $b\cdot v_1\sim
b\cdot v_2$, and $b\cdot v_i\ne0\cdot v_i$; so $gv_1\not\sim gv_2$,
and by Proposition~\ref{prop:onestep} the comparison is not an
equivalence for $n_L\ge3$. Proposition~\ref{prop:modular} does not
apply to this category: the only object with a non-trivial
automorphism is $2$, the only string with first object $2$ is $(2)$,
and for it $H=\Aut_{\cX}(2)$, so that $H\cap S_i=S_i$.
\end{ex}

The category of Example~\ref{ex:top} is the first of a family in which
the condition of Theorem~\ref{thm:main} can be read off the lattice of
subgroups of an abelian group.

\begin{prop}[a three-level family]\label{prop:dist}
Let $G$ be a finite abelian group and $A,T_1,\dots,T_m$ subgroups of
$G$. Let $\cX$ have the objects $0_1,\dots,0_m$, $1$, $2$, with
$\Aut(2)=A$ and trivial automorphism groups otherwise, and the arrows
\[
\cX(1,2)=G,\qquad\cX(0_j,1)=\cX(0_j,2)=G/T_j\quad(1\le j\le m),
\]
composition being the addition of $G$. Let $n_L\ge3$. Then $c^L_\cP$
is an equivalence if and only if
\begin{equation}\label{eq:dist}
A\cap\sum_{j\in J}T_j=\sum_{j\in J}(A\cap T_j)\qquad\text{for every }
J\subseteq\{1,\dots,m\}.
\end{equation}
\end{prop}
\begin{proof}
$\cX$ is an EI-category, and its strings have length at most three, so
that $\cS_n=\cS_3$. Choose representatives $\gamma$ of the cosets of
$A$ in $G$, with $0$ for $A$. Up to isomorphism the strings with last object $2$ are
$s=(2)$, $t_\gamma=(1\xrightarrow{\gamma}2)$, $u_{j,x,\gamma}=(0_j
\xrightarrow{x}1\xrightarrow{\gamma}2)$ with $x\in G/T_j$, and strings
$p=(0_j\to2)$. The arrows of $\cS_3$ between them are, besides the
identities and the automorphisms $a\in A$ of $s$: the arrows of $\cS_3
(t_\gamma,s)=A+\gamma\subseteq G$; the arrow $x\colon u_{j,x,\gamma}\to
t_\gamma$; the arrows of $\cS_3(u_{j,x,\gamma},s)=A+\gamma+x\subseteq
G/T_j$; and arrows $p\to s$. An arrow of $\cS_3$ preserves the last
object of a string.

Let $s'$ be a string other than $(2)$; its first object is $1$ or some
$0_j$. If $v\colon t'\to t$ and $u\colon t'\to s'$ are arrows of
$\cS_3$ and $g\colon t_0\to s'_0$ an arrow of $\cX$ with $gv=u$, then
$t_0\ne2$, and either $t_0=1$, in which case $s'_0=1$, the strings
$t$ and $s'$ are the final segment of $t'$ beginning at $1$, and
$g=\id$; or $t_0=0_j$, in which case $t=t'$, $v=\id$ and $g=u$. So the argument of
Proposition~\ref{prop:T} shows that the unit is injective at $s'$.

Let now $s=(2)$, and $W\colon\cS_3\to\Fin$. For $v\in W(t_\gamma)$ let
$J(v)$ be the set of the $j$ such that $v=W(x)y$ for some $x\in G/T_j$
and $y\in W(u_{j,x,\gamma})$, let $T(v)=\sum_{j\in J(v)}T_j$, and
$f_v\colon A+\gamma\to W(s)$, $f_v(\kappa)=W(\kappa)v$. Two remarks:
\begin{itemize}
\item[(a)] $f_v(\kappa+a)=f_v(\kappa)$ for $a\in A\cap T_j$, $j\in
J(v)$: if $v=W(x)y$, then $f_v(\kappa)=W(\kappa+x)y$, and $\kappa+a+x=
\kappa+x$ in $G/T_j$;
\item[(b)] in the colimit of $W$ over $\cJ_s$, the element $v$ at
$(t_\gamma,g)$ equals $v$ at $(t_\gamma,g+\tau)$, for $g\in G$, $\tau
\in T_j$, $j\in J(v)$: both are equal to $y$ at $(u_{j,x,\gamma},g+x)$.
\end{itemize}
By (b), for $\kappa,\kappa'\in A+\gamma$ with $\kappa-\kappa'\in T(v)$
the unit sends $f_v(\kappa)$ and $f_v(\kappa')$ to the same element,
since $f_v(\kappa)$ at $(s,\id)$ equals $v$ at $(t_\gamma,\kappa)$.

\emph{Necessity.} Suppose that \eqref{eq:dist} fails for $J$, and let
$\kappa\in A\cap\sum_{j\in J}T_j$ be outside $\sum_{j\in J}(A\cap
T_j)$. Let $u_j=u_{j,0,0}$, and let $W$ be the colimit of the
morphisms $\cS_3(t_0,-)\to\cS_3(u_j,-)$, $j\in J$, given by
composition with $0\colon u_j\to t_0$; so $W$ is generated by elements
$y_j\in W(u_j)$ with $W(0)y_j=v$ for all $j\in J$, $v$ being the class
of the identity of $t_0$. The set $\cS_3(u_j,s)$ is the image of $A$ in
$G/T_j$, so $W(s)$ is the colimit of the quotient maps $A\to A/(A\cap
T_j)$, $j\in J$, which is $A/\sum_{j\in J}(A\cap T_j)$, and $f_v$ is
the quotient map. So $f_v(\kappa)\ne f_v(0)$; since $J\subseteq J(v)$,
these elements have the same image under the unit.

\emph{Sufficiency.} Assume \eqref{eq:dist}, and that $W(s)$ is not
empty, with an element $w_0$. Let $\cJ^*$ be the full subcategory of
$\cJ_s$ of the objects $(t,g)$ with $t$ one of $s$, $t_\gamma$,
$u_{j,x,\gamma}$, or with $t=p$ and $g\in\cS_3(p,s)$; it is a union of
connected components of $\cJ_s$: an arrow of $\cS_3$ preserves the
last object, and an object $(p,c)$ with $c\notin\cS_3(p,s)$ has no
arrows to or from other objects, as $c=av$ with $a\in A$ and $v\in
\cS_3(p,s)$ is in $\cS_3(p,s)$. We
define a map $\rho$ from the colimit of $W$ over $\cJ^*$ to $W(s)$.
At $(s,a)$ and at $(p,c)$ let $\rho(w)=W(a)w$ and $\rho(w)=W(c)w$. At
$(t_\gamma,g)$ let $\rho(v)=f_v(\kappa)$ if there is $\kappa\in
(g+T(v))\cap(A+\gamma)$, and $\rho(v)=w_0$ otherwise; two such
$\kappa$ differ by an element of $A\cap T(v)$, which by
\eqref{eq:dist} is the sum of the $A\cap T_j$, $j\in J(v)$, so by (a)
this does not depend on $\kappa$. At $(u_{j,x,\gamma},c)$ let $\rho(y)$
be the value of $\rho$ on $W(x)y$ at $(t_\gamma,g)$, for $g\in G$ with
$g+x=c$; since $g$ is determined modulo $T_j$ and $j\in J(W(x)y)$,
this does not depend on $g$. The map $\rho$ is compatible with the
arrows of $\cJ^*$: with $x\colon(u_{j,x,\gamma},g+x)\to(t_\gamma,g)$ by
definition; with $\kappa\colon(t_\gamma,a+\kappa)\to(s,a)$ since
$\rho(v)=f_v(a+\kappa)=W(a)W(\kappa)v$ at $(t_\gamma,a+\kappa)$; with
$\kappa+x\colon(u_{j,x,\gamma},a+\kappa+x)\to(s,a)$, $\kappa\in
A+\gamma$, since $\rho(y)=f_{W(x)y}(a+\kappa)=W(a)W(\kappa+x)y$; and
with the arrows into $(s,a)$ from the objects $(s,a')$ and $(p,c)$ by
the functoriality of $W$. Since $\rho(w)=w$ at $(s,\id)$, the unit is
injective at $s$.
\end{proof}

\begin{ex}\label{ex:dist}
(a) Example~\ref{ex:top} is the case $G=\mathbb Z/2\times\mathbb Z/2$,
$A=\langle b\rangle$, $T_1=\langle a\rangle$, $T_2=\langle a+b
\rangle$: here $A\cap(T_1+T_2)=A$ and $A\cap T_j=0$.

(b) \emph{The retractions of Proposition~\ref{prop:retract} are not
necessary.} Let $G=(\mathbb Z/2)^3$ with basis $e_1,e_2,e_3$, $A=
\langle e_1,e_2\rangle$ and
\[
T_1=\langle e_1,e_3\rangle,\qquad T_2=\langle e_2,e_1+e_3\rangle,
\qquad T_3=\langle e_1+e_2,e_3\rangle,
\]
a category with $28$ arrows. The groups $A\cap T_j$ are $\langle e_1
\rangle$, $\langle e_2\rangle$, $\langle e_1+e_2\rangle$, and any two
of them generate $A$; so \eqref{eq:dist} holds and $c^L_\cP$ is an
equivalence for every $L$ proving linearity. For $s=(2)$ and $t=(1
\xrightarrow{0}2)$ the arrow $e_3\colon1\to2$ is in $C_s(t)$ and not in
$\cS_3(t,s)=A$: for the arrow $x\colon(0_j\xrightarrow{x}1
\xrightarrow{0}2)\to t$ the composite $e_3+x$ is an arrow into $s$ of
$\cS_3$, as $A+T_j=G$. A retraction would send $e_3$ to $\kappa\in A$
with $\kappa+x=e_3+x$ in $G/T_j$ for $j=1,2,3$, that is with $\kappa$
in each of the sets $(e_3+T_j)\cap A$, which are $\{0,e_1\}$,
$\{e_1,e_1+e_2\}$ and $\{0,e_1+e_2\}$; these meet in pairs and have no
common element. So for $n\ge3$ there is no retraction as in
Proposition~\ref{prop:retract}.

(c) \emph{The obstruction of Proposition~\ref{prop:onestep} is not the
only one.} Let $m\ge3$, $G=(\mathbb Z/2)^m$ with basis $e_1,\dots,e_m$,
$T_j=\langle e_j\rangle$ and $A=\langle e_1+\dots+e_m\rangle$; for
$m=3$ this is a category with $38$ arrows. Then $A\cap\sum_{j\in J}T_j$
is $A$ for $J=\{1,\dots,m\}$ and $0$ otherwise, and $A\cap T_j=0$; so
$c^L_\cP$ is not an equivalence for $n_L\ge3$, while it is an
equivalence for the full subcategory obtained by omitting any one of
the objects $0_j$. The hypothesis of Proposition~\ref{prop:onestep} is
never satisfied: it requires $(t,g)$ in $\cJ^c_s$ outside $\cJ^0_s$,
so, by the proof of Proposition~\ref{prop:dist}, $s=(2)$ and, an
object $u_{j,x,\gamma}$ receiving no arrow but the identity, $t=
t_\gamma$, $g\in G$ outside $A+\gamma$, and, $gv_i$ being an arrow of
$\cS_3$, $v_i\ne\id$, that is $v_i=x_i\colon
u_{j_i,x_i,\gamma}\to t_\gamma$, with $g\in A+\gamma+T_{j_i}$ as
$gv_i$ is an arrow of $\cS_3$. If $j_1\ne j_2$, then $g-\gamma\in
(A+T_{j_1})\cap(A+T_{j_2})=A$, which is excluded. If $j_1=j_2$, then
$g=\kappa+\tau$ with $\kappa\in A+\gamma$ and $\tau\in T_{j_1}$, so
that $gv_i=\kappa v_i$ and $gv_1\sim gv_2$.
\end{ex}

In general the condition of Theorem~\ref{thm:main} has to be tested on
finitely many functors only.

\begin{prop}[a finite test]\label{prop:test}
For a string $t$ of length at most $n$ let $I_t\colon\cS_n\to\Fin$ be
the functor sending $t'$ to the set of the subsets of $\cS_n(t',t)$,
and $k\colon t'\to t''$ to the map $U\mapsto\{l\in\cS_n(t'',t)\mid
lk\in U\}$. Then $c^L_\cP$ is an equivalence if and only if the unit
$I_t\to\pi^*\mathrm{Lan}_\pi I_t$ is injective for every $t$. In
particular it is decidable, for a finite EI-category $\cX$ and $n\in
\{1,2,\dots,\infty\}$, whether $c^L_\cP$ is an equivalence for the
logics $L$ with $n_L=n$.
\end{prop}
\begin{proof}
For $W\colon\cS_n\to\Fin$ the morphisms $\theta\colon W\to I_t$
correspond to the subsets $B$ of $W(t)$, by $B=\{w\mid\id_t\in
\theta_t(w)\}$ and $\theta_{t'}(w)=\{l\mid W(l)w\in B\}$. For $w\ne w'$
in $W(t)$ the morphism given by $B=\{w\}$ separates $w$ and $w'$; so
$W$ is a subfunctor of a finite product of functors $I_t$. The unit of
$W$ is injective if and only if $W$ is a subfunctor of $\pi^*Y$ for
some $Y\colon\cX\to\Fin$, every morphism $W\to\pi^*Y$ factoring through
the unit; and the functors with this property are closed under
subfunctors and, $\pi^*$ preserving products, under finite products. So
the units of all $W$ are injective as soon as those of the $I_t$ are,
and Theorem~\ref{thm:main} applies. The sets $I_t(t')$ and the values
of $\mathrm{Lan}_\pi I_t$ are finite and computable, $\cS_n$ being a
finite category.
\end{proof}

\section{Questions}

\begin{que}
Give a closed description of the finite EI-categories $\cX$ for which
every finite functor on $\cS_n(\cX)$ embeds in one pulled back from
$\cX$. By Proposition~\ref{prop:test} the condition is decidable; by
Example~\ref{ex:dist} the retractions of Proposition~\ref{prop:retract}
are not necessary for it, and the obstructions of
Propositions~\ref{prop:modular} and~\ref{prop:onestep} are not the
only ones. Is there a description in terms of the lattices of
subgroups of the automorphism groups of $\cX$, of which
Proposition~\ref{prop:modular} and the equalities \eqref{eq:dist} are
instances?
\end{que}

\begin{que}
Describe the models of $\HH_L\cP$ when $\cX$ is a finite category
with split idempotents which is not an EI-category.
Lemma~\ref{lem:linear} still holds, but a non-invertible endomorphism
has a power which is a non-identity idempotent, whose splitting is a
pair of non-invertible arrows with identity composite, so that there
are strings of unbounded length. A related example outside the present
setting, as it has infinite models: for the theory of pure equality
the comparison is an equivalence only for classical logic
\cite{Interm}.
\end{que}

\begin{que}
For the logics which prove the weak excluded middle and not linearity
the comparison fails already for the four-element chain \cite{Interm}.
Is there a description of the models of the translated theory over a
finite poset, in which the chains are replaced by the finite rooted
frames of the logic?
\end{que}

\section*{Declarations}
The author has no competing interests to declare.

\begin{thebibliography}{99}
\bibitem{AGM08} S.~Aguzzoli, B.~Gerla, V.~Marra, \emph{G\"odel algebras
free over finite distributive lattices}, Ann. Pure Appl. Logic 155
(2008), 183--193.
\bibitem{Carai} L.~Carai, \emph{Free algebras and coproducts in
varieties of G\"odel algebras}, preprint, arXiv:2406.05480.
\bibitem{Horn69} A.~Horn, \emph{Logic with truth values in a linearly
ordered Heyting algebra}, J. Symbolic Logic 34 (1969), 395--408.
\bibitem{Elephant} P.~T.~Johnstone, \emph{Sketches of an Elephant: A Topos
Theory Compendium}, Oxford Univ. Press, 2002.
\bibitem{Heyting} R.~Shibusawa, \emph{The free Heyting category on a
pretopos is a pretopos only in the Boolean case: the intuitionistic
translation and the pretopos completion}, preprint (2026), archived at
DOI 10.5281/zenodo.24182122.
\bibitem{Interm} R.~Shibusawa, \emph{Intermediate logics and
imaginaries: the translation into an intermediate logic commutes with
the pretopos completion of all finite propositional theories exactly
for the extensions of G\"odel--Dummett logic}, preprint (2026),
archived at DOI 10.5281/zenodo.24150610.
\end{thebibliography}

\end{document}
